\documentclass[11pt,a4paper]{article} % use larger type; default would be 10pt
\NeedsTeXFormat{LaTeX2e}
\makeindex

%---------------------------- PACKAGE INCLUSION -------------------------------
% This group renders characters clearer and more precise

\RequirePackage[bitstream-charter,cal,expert]{mathdesign}
\RequirePackage{charter}
\RequirePackage{helvet}
\RequirePackage{makeidx}
\RequirePackage{latexsym}

\usepackage{geometry}
\geometry{a4paper,
top=3.5em,
left=3em,
right=3em,
bottom=3.39em
}

\usepackage{graphicx}
\usepackage{adjustbox}
\usepackage{xspace}
\usepackage[parfill]{parskip} % Activate to begin paragraphs with an empty line rather than an indent
\usepackage{paralist} % very flexible & customisable lists (eg. enumerate/itemize, etc.)
\usepackage{listings} % for lstset definitions
\usepackage{url}
\usepackage{subfig} % make it possible to include more than one captioned figure/table in a single float
\usepackage{epsfig}
\usepackage{gensymb}
\usepackage{textcomp}
\usepackage{booktabs}

\usepackage{stackengine}
\newcommand\xrowht[2][0]{\addstackgap[.5\dimexpr#2\relax]{\vphantom{#1}}}

\usepackage{amsmath}

\usepackage{hyperref}
\hypersetup{
colorlinks,
pagebackref,
citecolor=medgreen,
linkcolor=purplish,
breaklinks,
pdftex,
bookmarks,
plainpages=false,
pdftitle={Line Plot representing "YERITHLINEPLOTTITLE"
from YERITHVENTESDEBUT to YERITHVENTESFIN (YERITHENTREPRISE)},
pdfauthor={YERITHENTREPRISE (YERITHUTILISATEUR)}
}


\usepackage{dataplot}


\definecolor{forestgreen}{RGB}{2,160,70}
\definecolor{mediumblue}{RGB}{27,93,255}
\definecolor{firebrickred}{RGB}{178,34,34}
\definecolor{listingray}{gray}{0.9}
\definecolor{lbcolor}{rgb}{0.9,0.9,0.9}
\definecolor{darkgreen}{rgb}{0,0.35,0}
\definecolor{medgreen}{rgb}{0,0.5,0}
\definecolor{lightgreen}{rgb}{0.5,0.7,0.5}
\definecolor{medgrey}{rgb}{0.6,0.6,0.6}
\definecolor{purplish}{rgb}{0.4,0,0.6}
\definecolor{brightred}{rgb}{1,0.2,0.2}


YERITHDTLSETBARCOLOR


% format two pieces of text, one left aligned and one right aligned
\newcommand{\headerrow}[2]
{\begin{tabular*}{\linewidth}{l@{\extracolsep{\fill}}r}
#1 &
#2 \\
\end{tabular*}}

\newcommand{\emphbold}[1]{\textbf{\emph{#1}}\xspace}


\begin{document}

\bigskip

\headerrow
	{\emphbold{YERITHENTREPRISE}}
	{\emph{\textbf{Tax payer registration n\textsuperscript{o}:} YERITHCONTRIBUABLENR}}
\headerrow
	{\emphbold{YERITHACTIVITESENTREPRISE}}
	{\emph{\textbf{Bank account n\textsuperscript{o}:} YERITHCOMPTEBANCAIRENR,}}
\headerrow
	{\emphbold{Email: YERITHEMAIL}}
	{\emph{YERITHAGENCECOMPTEBANCAIRE}}
\headerrow
	{\emphbold{Phone: YERITHTELEPHONE}}
	{}
\headerrow
	{\emphbold{P.O Box: YERITHBOITEPOSTALE YERITHVILLE}}
	{}

\hrule

\headerrow
	{}
	{\textbf{YERITHDATE}}


\section*{Line plot representing "YERITHLINEPLOTTITLE".}
\textbf{Generated by:} YERITHUTILISATEUR\\
\textbf{FILE (this) GENERATION Site:} YERITHSUCCURSALE\\
\textbf{Generated at:} YERITHHEUREGENERATION\\
\textbf{Time period:} YERITHVENTESDEBUT -- YERITHVENTESFIN


\begin{figure}[!htbp]
\setlength{\DTLbaroutlinewidth}{1pt}
\centering
%\DTLloaddb{data1}{data1.csv}
YERITHLINEPLOTDATA
\DTLplot{data1,data2}{x=Total,y=Nom,legend,box,style=lines}
\captionof{figure}{YERITHLINEPLOTTITLE}
\end{figure}


YERITHLINEPLOTFIN

\end{document}

